% supplementary_material_twistor.tex
%
% Build with:  make twistor-paper     (from the repository root)
% or:          python3 paper/make_strings_facts.py --sections twistor
%              && pdflatex supplementary_material_twistor.tex (x2)
%
% Every number quoted here is recomputed from the pyCICY source by
% paper/make_strings_facts.py and written to figures/twistor_facts.tex, which
% is \input{} below. The prose and the code cannot drift apart silently.

\documentclass[11pt,a4paper]{article}

\usepackage[T1]{fontenc}
\usepackage[utf8]{inputenc}
\usepackage[margin=1in]{geometry}
\usepackage[protrusion=true,expansion=false]{microtype}
\usepackage{amsmath,amssymb,amsthm}
\usepackage{graphicx}
\usepackage{booktabs}
\usepackage{enumitem}
\usepackage{xcolor}
\usepackage[hidelinks]{hyperref}
\usepackage[nameinlink,capitalize]{cleveref}
\usepackage{framed}
\usepackage{orcidlink}

\graphicspath{{figures/}}

% Numbers computed by paper/make_strings_facts.py --sections twistor.
% Declared here as fallbacks so the document typesets before the script has
% run; the \input below overrides every one with the computed value.
\newcommand{\TFBcfwMax}{?}\newcommand{\TFBcfwAgree}{?}
\newcommand{\TFBcfwAntiAgree}{?}\newcommand{\TFSixPointMHV}{?}
\newcommand{\TFSixPointBCFW}{?}\newcommand{\TFThreePtSquares}{?}
\newcommand{\TFThreePtAngles}{?}\newcommand{\TFThreePtHol}{?}
\newcommand{\TFThreePtAnti}{?}\newcommand{\TFNmhvValue}{?}
\newcommand{\TFNmhvCyclic}{?}\newcommand{\TFNmhvReflect}{?}
\newcommand{\TFRelationMax}{?}\newcommand{\TFUoneZero}{?}
\newcommand{\TFBcjZero}{?}\newcommand{\TFRankFive}{?}
\newcommand{\TFRankSix}{?}\newcommand{\TFOrderingsFive}{?}
\newcommand{\TFOrderingsSix}{?}\newcommand{\TFKKFive}{?}
\newcommand{\TFKKSix}{?}\newcommand{\TFBCJcountFive}{?}
\newcommand{\TFBCJcountSix}{?}\newcommand{\TFRankAgreesFive}{?}
\newcommand{\TFRankAgreesSix}{?}\newcommand{\TFCellsthree}{?}
\newcommand{\TFCellsfour}{?}\newcommand{\TFCellsfive}{?}
\newcommand{\TFCellTotalthree}{?}\newcommand{\TFCellTotalfour}{?}
\newcommand{\TFCellTotalfive}{?}\newcommand{\TFCellClosedthree}{?}
\newcommand{\TFCellClosedfour}{?}\newcommand{\TFCellClosedfive}{?}
\newcommand{\TFCellsGTwoFour}{?}\newcommand{\TFCellsSymmetric}{?}
\newcommand{\TFTopCellGTwoFour}{?}\newcommand{\TFChiPT}{?}
\newcommand{\TFChiIncidence}{?}\newcommand{\TFChiMinkowski}{?}
\newcommand{\TFKleinConfig}{?}\newcommand{\TFDegreeMHV}{?}
\newcommand{\TFDegreeNMHV}{?}\newcommand{\TFDegreeNMHVLoop}{?}
\newcommand{\TFPenroseGraviton}{?}\newcommand{\TFPenrosePhoton}{?}
\newcommand{\TFPenroseScalar}{?}\newcommand{\TFPenroseHOne}{?}
\IfFileExists{figures/twistor_facts.tex}{\input{figures/twistor_facts.tex}}{}

\newcommand{\PP}{\mathbb{P}}
\newcommand{\cO}{\mathcal{O}}
\newcommand{\QQ}{\mathbb{Q}}
\newcommand{\ket}[1]{\langle #1 \rangle}

\theoremstyle{definition}
\newtheorem{remark}{Remark}

\title{\vspace{-1cm}\textbf{What Holomorphy Determines,\\
and What a Rank Cannot See:\\[2pt]
\large Exact Methods for Twistor Geometry and Tree Amplitudes}}
\author{B.S. Hartshorn \orcidlink{0009-0004-2853-655X} \small \url{https://github.com/brentharts/CICY}}

\begin{document}
\maketitle

\begin{abstract}
A tree-level scattering amplitude is a rational function of spinor brackets, 
so with rational spinors it is a rational number, and every assertion about it is decidable. 
We construct momentum-conserving kinematics over $\QQ$ --- 
momentum conservation is a linear condition once the two spinors are
taken independent, so a configuration is built rather than solved for --- and
use it to check the Britto--Cachazo--Feng--Witten recursion against the
Parke--Taylor formula as an exact equality of rationals, for all
$n \leq \TFBcfwMax$ and for the conjugate amplitudes as well. Beyond the
maximally-helicity-violating sector there is no closed form to compare
against, and the amplitudes are instead pinned by symmetries that the construction does not manifest.

The $(n-1)!$ colour orderings of a tree amplitude
obey two families of linear relations. Computing the rank of the matrix of all
orderings evaluated at many kinematic points returns $(n-2)!$, the
Kleiss--Kuijf count, and never $(n-3)!$, the Bern--Carrasco--Johansson count
--- although the BCJ relation itself is verified to vanish exactly at every point tested. 
The reason is structural: a rank taken across kinematic points
can only detect relations whose coefficients are constant, and the BCJ
coefficients are Mandelstam invariants. The blind spot is a property of the
method rather than a failure of the physics, and stating it is more useful than engineering around it. 
Finally, the twistor double fibration is computed:
compactified complexified Minkowski space is the Klein quadric, a
configuration matrix, and its Euler characteristic
$\TFChiMinkowski$ counts the Schubert cells of $G(2,4)$ without anything being supplied.
\end{abstract}

\tableofcontents

% =====================================================================

\section{Introduction}
\label{sec:intro}

The first paper in this series~\cite{SuppMath} asked which quantities in a
heterotic compactification follow from the configuration matrix alone; the
second~\cite{SuppStrings} asked the same of F-theory and Type IIB
orientifolds, where the determining structure turned out to be anomaly
cancellation rather than geometry. Both were organized around a boundary: the
quantities that are exact, and the several distinguishable ways in which the
rest fail to be.

Scattering amplitudes sit differently. There is no compactification, no
metric to be obstructed by, and --- at tree level --- no integral. Penrose's
twistor correspondence replaces a point of spacetime by a line in
$\PP^3$~\cite{Penrose1967,Penrose1969}, and Witten's observation is that the
perturbative expansion of gauge theory, which in Feynman diagrams grows
faster than exponentially, becomes a statement about holomorphic curves in
that $\PP^3$~\cite{Witten2004}. The consequence exploited here is narrower and
more mundane than the twistor-string programme it launched: an $n$-gluon tree
amplitude is a rational function of the spinor brackets, so if the spinors are
rational the amplitude is a rational number.

That is what makes this paper possible in the same style as the earlier two.
Every equality below is an equality of elements of $\QQ$, computed with
\texttt{fractions.Fraction}, and every ``agrees'' is exact. There is no
floating point anywhere in \texttt{pyCICY.twistor}.

\Cref{sec:kinematics} sets up exact kinematics and explains why momentum
conservation is a linear condition rather than a constraint to be solved
numerically. \Cref{sec:pt,sec:bcfw} give the Parke--Taylor formula and check
the BCFW (Britto--Cachazo--Feng--Witten) recursion against it; \cref{sec:threepoint} isolates the three-point
degeneracy, which is a theorem rather than an inconvenience.
\Cref{sec:nmhv} goes beyond the closed forms. \Cref{sec:relations} is the
negative result. \Cref{sec:positroid,sec:geometry} treat the positive
Grassmannian and the twistor double fibration, the latter reusing the
complete-intersection machinery of the earlier papers unchanged.
\Cref{sec:boundary} restates the boundary for this setting, where it looks
different again.

% =====================================================================

\section{Exact kinematics}
\label{sec:kinematics}

A massless momentum in four dimensions is a rank-one bispinor,
$p^{a\dot b} = \lambda^a \tilde\lambda^{\dot b}$, and the Lorentz invariants
are the two brackets
%
\begin{equation}
\ket{ij} = \lambda_i^1 \lambda_j^2 - \lambda_i^2 \lambda_j^1 ,
\qquad
[ij] = \tilde\lambda_i^1 \tilde\lambda_j^2
     - \tilde\lambda_i^2 \tilde\lambda_j^1 .
\label{eq:brackets}
\end{equation}
%
Momentum conservation is
$\sum_i \lambda_i^a \tilde\lambda_i^{\dot b} = 0$.

Everything in this paper takes $\tilde\lambda$ \emph{independent} of
$\lambda$ rather than its complex conjugate. That is not a convenience. Real
Lorentzian three-point kinematics is degenerate --- see
\cref{sec:threepoint} --- so a recursion built from three-point amplitudes
exists only after complexification, and complexified kinematics with
independent spinors is precisely split signature. Over $\QQ$ it costs
nothing.

With the spinors independent, momentum conservation is linear. Choosing the
$\lambda_i$ freely, the condition says that each of the two columns of
$\tilde\lambda$ lies in the kernel of the $2 \times n$ matrix whose columns
are the $\lambda_i$. That kernel has dimension $n - 2$, and a rational vector
in it is computed by exact Gaussian elimination. A configuration is therefore
\emph{constructed} to conserve momentum, not fitted to; the method
\texttt{Kinematics.check} is a test of the implementation rather than of the
numbers, and it is one of the few checks in this paper that could not fail
for an interesting reason.

Two identities follow and neither is imposed anywhere. The Schouten identity
%
\begin{equation}
\ket{ij}\ket{kl} + \ket{ik}\ket{lj} + \ket{il}\ket{jk} = 0
\label{eq:schouten}
\end{equation}
%
holds because there is no antisymmetric three-index tensor in two dimensions;
it tests the bracket \cref{eq:brackets} rather than the kinematics. The
identity $\sum_i \ket{ki}[il] = 0$, which every amplitude manipulation uses,
is momentum conservation contracted with two spinors. Both come out exactly
zero, as does the Mandelstam invariant of the full set of momenta, and
$s_{ij}$ computed as the determinant of the summed bispinor agrees with
$\ket{ij}[ij]$ --- two routes to the same rational number.

Generic position matters and is enforced rather than hoped for: an amplitude
has poles wherever a consecutive bracket vanishes, so the constructor rejects
and redraws any configuration with a vanishing bracket rather than dividing by
zero later.

% =====================================================================

\section{Parke--Taylor}
\label{sec:pt}

For a colour-ordered tree amplitude with exactly two negative-helicity gluons
$i$ and $j$, Parke and Taylor's formula~\cite{ParkeTaylor1986}, proved by
Berends and Giele~\cite{BerendsGiele1988}, is
%
\begin{equation}
A(1, \dots, n) \;=\;
\frac{\ket{ij}^4}{\ket{12}\ket{23}\cdots\ket{n1}} .
\label{eq:pt}
\end{equation}
%
One line, replacing a number of Feynman diagrams that grows faster than
exponentially in $n$.

Three features of \cref{eq:pt} are checked rather than assumed. It is cyclic,
because the colour ordering is a cycle; it is invariant under reversal up to
$(-1)^n$; and its little-group weights are correct. Scaling $\lambda_i$ by
$t$ must scale the amplitude by $t^{-2h_i}$, so by $t^2$ on a
negative-helicity leg and $t^{-2}$ on a positive one. That is what fixes the
fourth power in the numerator: the denominator carries weight $-2$ on every
leg, so the numerator must supply $+4$ on each negative-helicity leg and
nothing anywhere else.

The formula depends on $\lambda$ alone. That is the geometric content, and is
Witten's starting point: an MHV amplitude \cite{CSW2004} is supported on a curve of degree
$\TFDegreeMHV$ in twistor space, which is to say on a line. On the fixed
kinematic point the module uses for six points, \cref{eq:pt} evaluates to
$\TFSixPointMHV$.

% =====================================================================

\section{BCFW, checked against the closed form}
\label{sec:bcfw}

The Britto--Cachazo--Feng--Witten recursion~\cite{BCF2005,BCFW2005} deforms
two of the external momenta by a complex parameter $z$,
%
\begin{equation}
\tilde\lambda_i \to \tilde\lambda_i - z \tilde\lambda_j ,
\qquad
\lambda_j \to \lambda_j + z \lambda_i ,
\label{eq:shift}
\end{equation}
%
which preserves both masslessness and momentum conservation, uses the fact
that the deformed amplitude vanishes as $z \to \infty$ for a suitable choice
of legs, and reconstructs the amplitude from its poles \cite{Hodges2009}. Each residue
factorizes into two amplitudes with fewer legs, so the whole tree expansion
follows from the three-point amplitudes, which Lorentz invariance alone fixes.

Over $\QQ$ this is entirely exact. The location $z_*$ of each pole is
rational, because the on-shell condition on the deformed internal momentum is
linear in $z$; the deformed internal momentum is then a rank-one bispinor and
factorizes into rational spinors. Nothing leaves the field at any stage.

The implementation shifts a cyclically adjacent pair with helicities
$(-,+)$, rotating the ordering to bring such a pair to the ends. Such a pair
always exists once both helicities are present, and rotating costs nothing
because the amplitude is cyclic.

\begin{framed}
\noindent
\textbf{The check.} BCFW and Parke--Taylor share no code: one is a recursion
seeded by three-point amplitudes, the other a closed formula. They agree
exactly for every $n$ up to $\TFBcfwMax$ (\TFBcfwAgree), and the conjugate
amplitudes --- $n-2$ negative helicities, obtained through the other
three-point amplitude --- agree with the parity-conjugate formula as well
(\TFBcfwAntiAgree). At the fixed six-point kinematic point of \cref{sec:pt},
the recursion returns $\TFSixPointBCFW$ against Parke--Taylor's
$\TFSixPointMHV$.
\end{framed}

\noindent
Writing this produced exactly one error, and the shape of the error is worth
recording because it is characteristic of sign conventions. An initial
version disagreed with \cref{eq:pt} by $(-1)^n$ for every $n$ tested. A
disagreement tracking the number of legs rather than the number of recursion
steps or the helicity configuration is the signature of a convention applied
once per amplitude, not of a mistake in any particular branch; it was traced
to a spurious sign on the anti-holomorphic three-point amplitude. A single
numerical evaluation would not have located it, and a floating-point
comparison would have shown a mismatch without showing the pattern.

The only input beyond the three-point amplitudes is that amplitudes with
fewer than two legs of either helicity vanish for $n \geq 4$. This is returned
rather than derived.

% =====================================================================

\section{Why three points need complex momenta}
\label{sec:threepoint}

At $n = 3$ the kernel that momentum conservation leaves is one-dimensional.
Both columns of $\tilde\lambda$ are then proportional to a single vector, so
every square bracket vanishes identically: on the holomorphic three-point
configuration the module builds, $\TFThreePtAngles$ of the three angle
brackets are non-zero and $\TFThreePtSquares$ of the three square brackets
are.

This is not an artifact of working over $\QQ$, and not a failure of the
generic-position constructor. It holds over any field: there is no three-point
configuration with both sets of brackets non-zero. The two three-point
amplitudes therefore live at \emph{different} kinematic points, and the module
exposes the choice explicitly rather than papering over it. On the
holomorphic point $A(1^-,2^-,3^+) = \TFThreePtHol$; on the anti-holomorphic
one $A(1^-,2^+,3^+) = \TFThreePtAnti$.

This degeneracy is exactly why BCFW requires complex momenta. With real
Lorentzian three-point kinematics every bracket vanishes and there is nothing
to recurse from. It is the one place in this paper where a computation
failing was the correct outcome, and the response was to make the failure a
documented feature rather than to route around it.

% =====================================================================

\section{Beyond the closed forms}
\label{sec:nmhv}

Three negative helicities at six points is the first amplitude with no closed
form, so there is nothing to compare a computation against. It is instead
pinned by symmetries that the construction does not make manifest.

The BCFW result is a sum over factorization channels selected \emph{after}
rotating the colour ordering to place a suitable shifted pair at the ends.
Nothing in that procedure makes the answer cyclic; each rotation produces a
genuinely different sum of genuinely different rational numbers. That the
totals agree is therefore a real check, and they do: the six-point NMHV
amplitude evaluates to $\TFNmhvValue$, is invariant under all six cyclic
rotations (\TFNmhvCyclic), and satisfies the reflection identity
(\TFNmhvReflect). The alternating-helicity configuration and a seven-point
$\mathrm{N}^2$MHV amplitude behave the same way.

Witten's degree formula records where these live. An
$\mathrm{N}^{k-2}\mathrm{MHV}$ amplitude at $\ell$ loops is supported on a
curve in twistor space of degree $d = k - 1 + \ell$ and genus at most $\ell$:
MHV on a line ($d = \TFDegreeMHV$), NMHV on a conic
($d = \TFDegreeNMHV$), and a loop raising the degree by one
($d = \TFDegreeNMHVLoop$ for NMHV at one loop). This is recorded rather than
derived; nothing else in the module uses it, and no loop amplitude is
computed anywhere.

% =====================================================================

\section{What a rank cannot see}
\label{sec:relations}

The $(n-1)!$ cyclically inequivalent colour orderings of a tree amplitude are
not independent. Two families of relations reduce them, and the two behave
completely differently under the obvious computational test.

The photon decoupling identity is the simplest. Summing over the cyclic
insertions of one leg into the ordering of the others gives zero:
%
\begin{equation}
A(1,2,\dots,n) + A(2,1,3,\dots,n) + \cdots
+ A(2,3,\dots,n-1,1,n) \;=\; 0 ,
\label{eq:u1}
\end{equation}
%
with $n-1$ terms rather than $n$, since inserting leg $1$ after leg $n$
returns the first ordering. The Kleiss--Kuijf relations~\cite{KleissKuijf1989}
generalise it, and all of them have coefficients $\pm 1$. The
Bern--Carrasco--Johansson relations~\cite{BCJ2008} go further; the
fundamental one is
%
\begin{equation}
\sum_{i=2}^{n-1} \Bigl( \sum_{j=2}^{i} s_{1j} \Bigr)\,
A(2, \dots, i, 1, i+1, \dots, n) \;=\; 0 ,
\label{eq:bcj}
\end{equation}
%
and its coefficients are Mandelstam invariants.

Both hold exactly. \Cref{eq:u1,eq:bcj} evaluate to zero for every $n$ up to
$\TFRelationMax$ tested, at every kinematic point
(\TFUoneZero{} and \TFBcjZero{} respectively). Counting independent orderings
should therefore be a matter of taking a rank: evaluate all $(n-1)!$ orderings
at many independent kinematic points and compute the rank of the resulting
matrix over $\QQ$.

\begin{center}
\begin{tabular}{lcccc}
\toprule
$n$ & orderings & rank & $(n-2)!$ & $(n-3)!$ \\
\midrule
$5$ & \TFOrderingsFive & \TFRankFive & \TFKKFive & \TFBCJcountFive \\
$6$ & \TFOrderingsSix  & \TFRankSix  & \TFKKSix  & \TFBCJcountSix  \\
\bottomrule
\end{tabular}
\end{center}

\noindent
The rank is $(n-2)!$, and it is never $(n-3)!$.

The explanation is structural rather than numerical, and it is the point of
this section. A rank taken across kinematic points detects only those linear
relations that hold \emph{simultaneously at every point}, which is to say only
relations with constant coefficients. Kleiss--Kuijf relations have
coefficients $\pm 1$, so they apply at every point at once and the rank
collapses to $(n-2)!$. The BCJ coefficients in \cref{eq:bcj} vary from point
to point, so no single linear combination of the columns annihilates the whole
matrix, and the further reduction is invisible --- even though, as the
residual shows, the relation holds exactly at each point separately.

Reporting $(n-3)!$ here would require supplying the relations as input, which
would make the computation a restatement of its own assumption rather than a
check on it. Reporting $(n-2)!$ and saying why is the informative outcome. It
is also a warning about a class of numerical experiment: a rank over sampled
parameter values measures the constant-coefficient part of a relation ideal
and nothing else, and an absent reduction is evidence about the coefficients,
not about the relations.

The implementation reports whether enough kinematic points were supplied for
the rank to have saturated, so a rank bounded by the sample size is visible
rather than mistaken for a result.

% =====================================================================

\section{The positive Grassmannian}
\label{sec:positroid}

Postnikov's classification~\cite{Postnikov2006} puts the cells of the totally
non-negative Grassmannian $G(k,n)_{\geq 0}$ in bijection with decorated
permutations of $[n]$ --- permutations with each fixed point coloured --- that
have $k$ anti-exceedances. The same objects label the on-shell diagrams of
planar $\mathcal{N}=4$ super Yang--Mills~\cite{ABCGPT2016, RSV2004}, and the
amplituhedron~\cite{ArkaniHamedTrnka2013} is built from them, which is why the
combinatorics belongs beside the amplitudes rather than in an appendix.

Counting directly:

\begin{center}
\begin{tabular}{lccc}
\toprule
$n$ & cells by $k$ & total & $\sum_j n!/j!$ \\
\midrule
$3$ & \TFCellsthree & \TFCellTotalthree
& \TFCellClosedthree \\
$4$ & \TFCellsfour & \TFCellTotalfour
& \TFCellClosedfour \\
$5$ & \TFCellsfive & \TFCellTotalfive
& \TFCellClosedfive \\
\bottomrule
\end{tabular}
\end{center}

\noindent
The rows are symmetric under $k \mapsto n-k$ (\TFCellsSymmetric), as duality
of Grassmannians requires; $G(2,4)$ has $\TFCellsGTwoFour$ cells and its top
cell has dimension $\TFTopCellGTwoFour = k(n-k)$.

The total is the check worth making. Summed over $k$ the count matches
$\sum_{j=0}^{n} n!/j!$, the number of arrangements of $n$ objects, and the two
sides have nothing obviously to do with one another: one counts permutations
with coloured fixed points graded by anti-exceedances, the other counts
ordered selections. The agreement tests the anti-exceedance statistic, which
is the part of the bijection easiest to get subtly wrong.

% =====================================================================

\section{The double fibration, as complete intersections}
\label{sec:geometry}

Twistor space is $\PP^3$. Compactified complexified Minkowski space is the
Grassmannian of two-planes in $\mathbb{C}^4$, which the Pl\"ucker embedding
realises as a single quadric in $\PP^5$ --- the Klein quadric. That is a
configuration matrix, $\TFKleinConfig$, so the complete-intersection machinery
of the earlier papers applies to it unchanged: the Euler characteristics below
come from the same Chern class routine written for orientifold fixed
loci~\cite{SuppStrings}.

\begin{center}
\begin{tabular}{lccc}
\toprule
space & configuration & $\dim$ & $\chi$ \\
\midrule
twistor space $PT$          & $[[3]]$              & 3 & \TFChiPT \\
incidence $F(1,3;4)$        & $[[3,1],[3,1]]$      & 5 & \TFChiIncidence \\
Minkowski $G(2,4)$          & $\TFKleinConfig$     & 4 & \TFChiMinkowski \\
\bottomrule
\end{tabular}
\end{center}

\noindent
The $\TFChiMinkowski$ is the number of Schubert cells of $G(2,4)$, indexed by
the partitions fitting in a two-by-two box, and the $\TFChiIncidence$ is the
cell count of the flag manifold, the $(1,1)$ hypersurface expressing the
incidence $Z^\alpha W_\alpha = 0$. Neither was supplied.

The Penrose transform~\cite{Penrose1969,PenroseMacCallum1972, ADM2017} sends a free
massless field of helicity $h$ to a class in $H^1$ of twistor space valued in
$\cO(-2h-2)$: $\TFPenroseScalar$ for a scalar, $\TFPenrosePhoton$ for a
photon, $\TFPenroseGraviton$ for a graviton. The twist is the content ---
helicity is the homogeneity degree of a function on twistor space, so the
entire spectrum of massless free fields is one sheaf cohomology with a varying
line bundle.

The caveat the module records is that the cohomology cannot be taken on
$\PP^3$. By Bott's formula $H^1(\PP^3, \cO(m)) = \TFPenroseHOne$ for every
$m$ \cite{Bott1957}, so the transform on compact twistor space would produce nothing at all.
It is taken on $\PP^3$ minus a line, or on a neighbourhood of a real slice,
and that open space is outside what this package's cohomology covers.

% =====================================================================

\section{The boundary}
\label{sec:boundary}

The two earlier papers each ended by classifying the ways a quantity can fail
to be available. Here the classification looks different again, and the
difference is informative about all three.

\begin{enumerate}[leftmargin=2em]
\item \emph{Exact and complete.} Every tree amplitude, for any number of legs
and any helicity assignment. There is no approximation, no metric and no
choice of data left open. This category is far larger here than in either
earlier paper, and the reason is that a tree amplitude is a rational function
and nothing else.

\item \emph{Exact but invisible to the method used.} The BCJ (Bern--Carrasco--Johansson) reduction of
\cref{sec:relations}. The relations are exact and are verified exactly; the
particular computation applied to them --- a rank over sampled kinematics ---
provably cannot see them. This is a category the earlier papers had no
occasion to name, and it is not the same as ``not implemented'': the obstacle
is a property of the question asked, not of the effort spent.

\item \emph{Not computed.} Loops. The tree expansion is exact and complete;
the loop expansion requires integrals over regions not determined by the
rational data here, and no loop amplitude is computed yet. The amplituhedron
appears only through the combinatorics of its cells, not its geometry.

\item \emph{On the wrong space.} The Penrose transform, which needs an open
twistor space that the package's cohomology machinery does not describe yet.
\end{enumerate}

\noindent
The second category is the one to carry forward. In~\cite{SuppMath} the
boundary was between what a configuration matrix contains and what it does
not; in~\cite{SuppStrings} between what consistency forces and what a choice
of flux or brane leaves open. Both are boundaries of the \emph{data}. Here the
data is complete --- the amplitudes are known exactly --- and the boundary is
of the \emph{method}. A rank computation is a perfectly good instrument that
measures one specific thing, and mistaking what it measures for what one
wanted to know is an error no amount of exactness protects against.

% =====================================================================

\section{Conclusion}
\label{sec:conclusion}

The first paper found that a configuration matrix determines more than one
expects \cite{SuppMath}. The second found that consistency does \cite{SuppStrings}. 
This one finds that holomorphy does: once an amplitude is known to be a rational function of
spinor brackets, working over $\QQ$ makes every question about it decidable,
and the resulting checks are equalities rather than agreements to some number of digits.

The methodological thread of the series holds here too. A quantity computed
once is a result; computed twice, by routes sharing no code, it is a test.
BCFW against Parke--Taylor, the cell counts against $\sum_j n!/j!$, the
Mandelstam invariant as a determinant against its bracket form, the Euler
characteristic of the Klein quadric against the Schubert cells of $G(2,4)$:
each pair is a check, and the one error made in writing this module was found
by such a pair disagreeing in a pattern --- $(-1)^n$ --- that identified the
convention responsible.

What is new here is the fourth item on that list of pairs never disagreeing
where it was expected to. The rank computation of \cref{sec:relations} does
not reproduce the BCJ count, and no amount of additional kinematic points will
make it. Exact arithmetic guarantees that a computed answer is the right
answer to the question posed; it says nothing about whether that was the
question one meant to ask.

% =====================================================================

\section{Source code and limitations}
\label{sec:source}
\footnotesize
Everything above is implemented in \texttt{pyCICY.twistor}, with
\texttt{tests/test\_twistor.py} executing the cross-checks quoted here and
\texttt{examples/twistor.py} reproducing the tables. The numbers in this
document are emitted by \texttt{paper/make\_strings\_facts.py} with
\texttt{--sections twistor}. The repository is at
\url{https://github.com/brentharts/CICY}.

Limitations:
%
\begin{itemize}[leftmargin=2em,itemsep=1pt]
\item Tree level only. No loop amplitude is computed, and the twistor-string
degree formula is recorded rather than derived or used.
\item Gluons only: colour-ordered pure Yang--Mills amplitudes, with no
supersymmetric multiplet structure, no fermions and no gravity.
\item The BCJ reduction is verified as an identity but not used to construct
a basis of orderings, for the reason given in \cref{sec:relations}.
\item The positive Grassmannian appears only combinatorially. Boundary
stratification, plabic graph moves \cite{Affolter} and the amplituhedron's geometry \cite{ArkaniHamedTrnka2013} are
absent.
\item The Penrose transform is recorded as a correspondence, not implemented;
the relevant cohomology is on an open twistor space.
\end{itemize}

% =====================================================================

\begin{thebibliography}{99}
\setlength{\itemsep}{1pt plus 0.5pt}

\bibitem{SuppMath}
B.S.~Hartshorn,
\emph{What a configuration matrix determines, and what it does not: exact
methods for heterotic compactifications on CICYs},
Pre-print (2026), \url{https://doi.org/10.5281/zenodo.21844007}.

\bibitem{SuppStrings}
B.S.~Hartshorn,
\emph{What anomaly cancellation determines, and what does not exist: exact
methods for F-theory and Type IIB orientifolds},
Pre-print (2026), \url{https://doi.org/10.5281/zenodo.21865768}.


\bibitem{Penrose1967}
R.~Penrose,
\emph{Twistor algebra},
J.\ Math.\ Phys.\ \textbf{8} (1967) 345.

\bibitem{Penrose1969}
R.~Penrose,
\emph{Solutions of the zero-rest-mass equations},
J.\ Math.\ Phys.\ \textbf{10} (1969) 38.

\bibitem{Witten2004}
E.~Witten,
\emph{Perturbative gauge theory as a string theory in twistor space},
Commun.\ Math.\ Phys.\ \textbf{252} (2004) 189,
\href{https://arxiv.org/abs/hep-th/0312171}{arXiv:hep-th/0312171}.

\bibitem{ParkeTaylor1986}
S.~J.~Parke and T.~R.~Taylor,
\emph{An amplitude for $n$-gluon scattering},
Phys.\ Rev.\ Lett.\ \textbf{56} (1986) 2459.

\bibitem{BerendsGiele1988}
F.~A.~Berends and W.~T.~Giele,
\emph{Recursive calculations for processes with $n$ gluons},
Nucl.\ Phys.\ B \textbf{306} (1988) 759. \url{https://doi.org/10.1016/0550-3213(88)90442-7}

\bibitem{CSW2004}
F.~Cachazo, P.~Svrcek and E.~Witten,
\emph{MHV vertices and tree amplitudes in gauge theory},
JHEP \textbf{09} (2004) 006,
\href{https://arxiv.org/abs/hep-th/0403047}{arXiv:hep-th/0403047}.

\bibitem{BCF2005}
R.~Britto, F.~Cachazo and B.~Feng,
\emph{New recursion relations for tree amplitudes of gluons},
Nucl.\ Phys.\ B \textbf{715} (2005) 499,
\href{https://arxiv.org/abs/hep-th/0412308}{arXiv:hep-th/0412308}.

\bibitem{BCFW2005}
R.~Britto, F.~Cachazo, B.~Feng and E.~Witten,
\emph{Direct proof of tree-level recursion relation in Yang--Mills theory},
Phys.\ Rev.\ Lett.\ \textbf{94} (2005) 181602,
\href{https://arxiv.org/abs/hep-th/0501052}{arXiv:hep-th/0501052}.

\bibitem{Hodges2009}
A.~Hodges,
\emph{Eliminating spurious poles from gauge-theoretic amplitudes},
JHEP \textbf{05} (2013) 135,
\href{https://arxiv.org/abs/0905.1473}{arXiv:0905.1473}.


\bibitem{KleissKuijf1989}
R.~Kleiss and H.~Kuijf,
\emph{Multigluon cross-sections and 5-jet production at hadron colliders},
Nucl.\ Phys.\ B \textbf{312} (1989) 616. \url{https://doi.org/10.1016/0550-3213(89)90574-9}

\bibitem{BCJ2008}
Z.~Bern, J.~J.~M.~Carrasco and H.~Johansson,
\emph{New relations for gauge-theory amplitudes},
Phys.\ Rev.\ D \textbf{78} (2008) 085011,
\href{https://arxiv.org/abs/0805.3993}{arXiv:0805.3993}.

\bibitem{Postnikov2006}
A.~Postnikov,
\emph{Total positivity, Grassmannians, and networks},
\href{https://arxiv.org/abs/math/0609764}{arXiv:math/0609764}.

\bibitem{ABCGPT2016}
N.~Arkani-Hamed, J.~Bourjaily, F.~Cachazo, A.~Goncharov, A.~Postnikov and
J.~Trnka,
\emph{Grassmannian Geometry of Scattering Amplitudes},
Cambridge University Press (2016),
\href{https://arxiv.org/abs/1212.5605}{arXiv:1212.5605}.

\bibitem{ArkaniHamedTrnka2013}
N.~Arkani-Hamed and J.~Trnka,
\emph{The amplituhedron},
JHEP \textbf{10} (2014) 030,
\href{https://arxiv.org/abs/1312.2007}{arXiv:1312.2007}.

\bibitem{RSV2004}
R.~Roiban, M.~Spradlin and A.~Volovich,
\emph{On the tree-level S-matrix of Yang--Mills theory},
Phys.\ Rev.\ D \textbf{70} (2004) 026009,
\href{https://arxiv.org/abs/hep-th/0403190}{arXiv:hep-th/0403190}.




\bibitem{PenroseMacCallum1972}
R.~Penrose and M.~A.~H.~MacCallum,
\emph{Twistor theory: an approach to the quantisation of fields and space-time},
Phys.\ Rept.\ \textbf{6} (1972) 241.

\bibitem{ADM2017}
M.~Atiyah, M.~Dunajski and L.~J.~Mason,
\emph{Twistor theory at fifty: from contour integrals to twistor strings},
Proc.\ Roy.\ Soc.\ A \textbf{473} (2017) 20170530,
\href{https://arxiv.org/abs/1704.07464}{arXiv:1704.07464}.


\bibitem{Bott1957}
R.~Bott,
\emph{Homogeneous vector bundles},
Ann.\ Math.\ \textbf{66} (1957) 203.

\bibitem{Affolter}
N.~Affolter, M.~Glick, P.~Pylyavskyy, S.~Ramassamy
\emph{Vector-relation configurations and plabic graphs} (version2 2023)
\url{https://arxiv.org/abs/1908.06959v2}

\end{thebibliography}

\end{document}
